\section{Literature Review}

In recent years, there has been a surge in research focused on the design and development of On-Board Diagnostic (OBD) systems, alongside the application of machine learning (ML) algorithms for predicting engine performance, optimizing fuel consumption, and detecting system anomalies. This section reviews key works in these domains and highlights existing research gaps relevant to the development of the ADAM system.

\subsection{Design and Development of On-Board Diagnostic Circuits}

The work presented in \cite{b1} outlines the implementation of an OBD system that collects real-time vehicular data such as engine speed, coolant temperature, and diagnostic trouble codes (DTCs) using a KEAZ128 microcontroller. The system emphasizes efficient data collection from the ECU and lays the groundwork for embedded diagnostic units.

In another study \cite{b2}, a multi-sensor data logger is developed, which includes GPS, an Inertial Measurement Unit (IMU), camera input, and OBD data acquisition. Data is logged locally every 100 ms using a BeagleBoard-XM platform powered by the TI OMAP3530 processor. This architecture allows for synchronized multi-modal vehicle monitoring, enabling rich data capture for analysis and event reconstruction.

The system proposed in \cite{b3} uses the ELM327 integrated circuit for interfacing with the OBD port, coupled with GPS data acquisition. Data is transmitted over Wi-Fi using the Carambola2 module and stored on a remote SQL server, highlighting the importance of cloud integration and remote diagnostics for modern vehicles.

\subsection{Review of Machine Learning Algorithms}

A comprehensive survey presented in \cite{b4} reviews a range of ML models used in vehicle diagnostics and engine control. It explores algorithms such as Artificial Neural Networks (ANN), Support Vector Machines (SVM), Random Forests (RF), and Adaptive Neuro-Fuzzy Inference Systems (ANFIS). The paper also notes the emerging relevance of Reinforcement Learning (RL) and Deep Reinforcement Learning (DRL) in optimizing fuel economy, controlling emissions, and enabling real-time adaptive control.

In the domain of network diagnostics, \cite{b5} introduces an ML-based fault diagnosis framework for Self-Organizing Networks (SONs). While this study is focused on mobile networks, it employs techniques such as Softmax-based Neural Networks and SVMs for multiclass classification, which can be analogously applied to vehicular diagnostic systems. Feature extraction is achieved using Key Performance Indicators (KPIs) and Performance Management Counters (PMCs), and the proposed solution is robust against both single and compound system faults.

Another study \cite{b6} investigates the use of feedforward multilayer perceptron ANNs to detect and classify internal combustion engine faults. The study compares various training algorithms—including backpropagation, Levenberg–Marquardt, quasi-Newton, and the newly proposed Smooth Variable Structure Filter (SVSF). Using crank-angle-converted vibration data from a V8 engine, the system successfully detected issues such as worn bearings and abnormal piston noise. Notably, the SVSF algorithm achieved an accuracy of 97\% and was able to detect and classify faults within 30 minutes, offering a cost-effective and fast solution for diagnostics.

In \cite{b7}, the authors present OBD-SecureAlert, a CAN-bus-based anomaly detection system designed to address the lack of security in in-vehicle networks. The system employs a Hidden Markov Model (HMM) to detect deviations in vehicle behavior by analyzing real-time parameters such as speed and engine RPM. Using a sliding window approach, it classifies sequences as either normal or anomalous based on predefined thresholds. Attack scenarios were simulated by injecting malicious data into the CAN bus, demonstrating the system's ability to detect cybersecurity threats effectively.

\textbf{Summary:} While significant progress has been made in both OBD system design and ML-based diagnostics, there remains a gap in integrating low-cost hardware, real-time edge analytics, secure telemetry, and intuitive user interfaces in a single cohesive system. The ADAM system addresses these limitations by combining embedded diagnostics, sensor fusion, machine learning, and mobile-based visualization in a lightweight and accessible architecture.


% \section{Literature Review:}
% In Recent years, there has been advancements in design and development of On board Diagnostic and also there are machine learning algorithms for predicting engine performance , fuel consumption and anomaly detection. In this section, we review the literature on design and development of OBD boards and developed machine learning algorithms.Additionally, we identify gaps in recent literature
% \subsection{Design and Development of On-Boards Diagnostic circuits}

% The study \cite{b1} contributes to the design of On Board Diagnostic systems that provides real-time information about a vehicle, such as engine speed, coolant temperature, pressure, and Diagnostic Trouble Codes (DTCs).Data Collection of OBD Using a keaz128 microcontroller.
% This paper \cite{b2} designs an Diagnostic system that has various data collection sources like gps,imu,camera and  OBD port and stores locally for every  100 ms. Using BeagleBoard-XM with TI OMAP3530 processor logs the data.
% This paper \cite{b3} designs of Diagnostic system that collect obd data using  ELM327 microcontroller along with gps .The data is transmitted to remote server and stored in sql database and transmit data using Carambola2 WiFi module.

% \subsection{Review of Machine Learning Algorithms}

% This study \cite{b4} is a detailed survey of machine learning methods used in engine performance, control, and fault diagnosis. It emphasizes the employment of Artificial Neural Networks (ANN), Support Vector Machines (SVM), Random Forests (RF), and neuro-fuzzy models such as ANFIS for improving fuel economy and emission control. The review also discusses the use of AI in engine control with a special focus on reinforcement learning (RL) and deep reinforcement learning (DRL) for enhanced adaptability, efficiency, and real-time decision-making in sophisticated engine operations.
% This paper \cite{b5} proposes an ML-based algorithm for automatic fault diagnosis in mobile networks using Self-Organizing Networks (SONs). It combines Softmax NN and SVM Models to carry out multiclass classification while diagnosing network faults. Along with KPIs and  performance management counters (PMCs) for feature extraction. This unsupervised approach handles both single and multiple faults, offering a robust and adaptable solution for modern mobile networks.
% This paper \cite{b6} examines the  use of artificial neural networks for the detection and classification of faults in internal combustion engines. It uses feedforward multilayer perceptron ANNs and training algorithms like backpropagation, Levenberg-Marquardt, quasi-Newton, extended Kalman filter, and a new Smooth Variable Structure Filter (SVSF). The study tested a method to detect engine faults without needing to know their exact location. Using vibration data from a V8 engine, converted into the crank angle domain, the system identified issues like missing bearings, chain tensioners and piston noise.Out of all the methods, SVSF, achieved 97\% accuracy and avoided common training errors. It can detect and classify faults in under 30 minutes, offering a fast, low-cost solution for engine diagnostics.
% This paper \cite{b7} presents OBD-SecureAlert, a data-driven system that detects anomalies in vehicle behavior using CAN bus data.Since CAN buses lack built-in security measures, it is open to cyberattacks that will compromise the safety of automobiles. To overcome this situation the system uses Hidden Markov Model (HMM) to monitor and detect deviations of normal vehicle behaviour, emphasizing on detection rather than prevention. It collects data such as Speed and RPM, trains it using a sliding window and flags the data sets based on predefined threshold limits as anomalies or safety/cybersecurity threats. Attack scenarios are simulated, like injecting anomalous data into CAN bus in order to detect anomalies.
% This paper \cite{b8} talks about DrivingStyles, a platform which aims to analyze and improve driving habits to diminish fuel consumption and greenhouse gases. It uses data from the ECU using an OBD-II Bluetooth interface, and smartphone technology to assess driving behavior and the efficiency of fuel. Applying data mining and neural networks, it classifies driver styles based on speed, acceleration, and engine RPM.
% This paper \cite{b9} talks about Enhanced wireless and computing technologies in vehicles driving the growth of mobile commerce (m-commerce) applications, such as in-car entertainment, real-time diagnostics, and connectivity services. These features offer greater convenience and smarter vehicle management for users.However, key challenges remain, including stable connectivity, driver safety, cybersecurity, and high deployment costs. 
